\documentclass[aps,prl,twocolumn,superscriptaddress]{revtex4-2}
\usepackage{amsmath,amssymb,amsfonts}
\usepackage{graphicx}
\usepackage{hyperref}

\begin{document}

\title{$U_q(\mathfrak{sl}_2)$ in Lean~4: The First Quantum Group\\
in a Proof Assistant}

\author{J.~Roehm}
\affiliation{Independent Researcher}

\date{\today}

\begin{abstract}
We present the first definition of a quantum group in any proof assistant:
the Drinfeld-Jimbo quantum group $U_q(\mathfrak{sl}_2)$, constructed in
Lean~4 with Mathlib as a quotient of the free algebra
$\mathrm{FreeAlgebra}(\mathbb{k}[q,q^{-1}], \{E,F,K,K^{-1}\})$ by the
Chevalley relations via \texttt{RingQuot} --- with zero axioms. The
\texttt{Ring} and \texttt{Algebra}\,$\mathbb{k}[q,q^{-1}]$ instances are
automatic. We formalize $q$-integers $[n]_q$ as Laurent polynomials with
the classical limit $[n]_1 = n$, the $q$-deformed Dolan-Grady coefficient
$[3]_q = q^2 + 1 + q^{-2}$, and the connection $[2]_1^4 = 16$ to the
Onsager algebra's DG coefficient. $K$ is proved invertible with $K^{-1}$
as its inverse, and all five Chevalley relations are stated as theorems.
This establishes the algebraic foundation for the
$q$-Onsager $\to$ quantum group $\to$ gauge emergence chain, connecting
lattice chirality (formalized in companion work) to continuous gauge theory.
The construction uses Mathlib's existing \texttt{FreeAlgebra},
\texttt{RingQuot}, and \texttt{LaurentPolynomial} infrastructure,
following the same pattern as \texttt{TensorAlgebra} and
\texttt{UniversalEnvelopingAlgebra}.
\end{abstract}

\maketitle

%% ══════════════════════════════════════════════════════════════════
\section{Introduction}

Quantum groups, introduced independently by Drinfeld~\cite{Drinfeld1986}
and Jimbo~\cite{Jimbo1986}, are $q$-deformations of universal enveloping
algebras that play a central role in knot theory, integrable systems,
topological quantum field theory, and the mathematical foundations of
gauge theory. Despite their importance and 40~years of development, no
quantum group has previously been defined in a proof assistant.

We present the first such definition: $U_q(\mathfrak{sl}_2)$ in Lean~4
with Mathlib, constructed from \texttt{FreeAlgebra} and \texttt{RingQuot}
with \emph{zero axioms}. The key insight is that Mathlib's existing
algebra-by-generators-and-relations infrastructure --- the same
pattern used for \texttt{TensorAlgebra}, \texttt{ExteriorAlgebra},
and \texttt{UniversalEnvelopingAlgebra} --- directly applies to quantum
groups. Combined with Mathlib's \texttt{LaurentPolynomial} type
(providing the base ring $\mathbb{k}[q, q^{-1}]$) and the recently merged
\texttt{Coalgebra} $\to$ \texttt{Bialgebra} $\to$ \texttt{HopfAlgebra}
hierarchy, the infrastructure for quantum group formalization is complete.

This work is part of a broader program formalizing the mathematical
foundations of gauge emergence from topological order. In companion
work~\cite{Paper9,Paper10}, we formalized the Onsager algebra (24
theorems), the Drinfeld center $Z(\mathrm{Vec}_G) \cong \mathrm{Rep}(D(G))$
(87 theorems), and the generation constraint $N_f \equiv 0 \pmod{3}$
from modular invariance. The quantum group $U_q(\mathfrak{sl}_2)$ is
the next link in the chain: at $q = 1$, it recovers the classical
$U(\mathfrak{sl}_2)$ and the Onsager algebra's Chevalley embedding;
at roots of unity $q = e^{2\pi i/(k+2)}$, its representation category
becomes a modular tensor category equivalent to $SU(2)_k$
Chern-Simons theory.

%% ══════════════════════════════════════════════════════════════════
\section{$q$-Integers as Laurent Polynomials}

The $q$-integer is defined as a sum of Laurent monomials:
\begin{equation}
[n]_q = \sum_{i=0}^{n-1} q^{n-1-2i} = q^{n-1} + q^{n-3} + \cdots + q^{-(n-1)},
\end{equation}
a sum of $n$ terms. This \emph{sum form} avoids the fraction
$(q^n - q^{-n})/(q - q^{-1})$ and is well-defined for all $q$,
including $q = 1$. In Lean~4, we define \texttt{qInt n} as an element
of \texttt{LaurentPolynomial k}, where $T^m$ represents $q^m$.

Explicit values:
\begin{align}
[0]_q &= 0, \quad [1]_q = 1, \quad [2]_q = q + q^{-1}, \\
[3]_q &= q^2 + 1 + q^{-2}.
\end{align}
The \emph{classical limit} $[n]_1 = n$ holds because each $q^m$
evaluates to $1$ at $q = 1$, giving a sum of $n$ ones.

The coefficient $[3]_q$ appears in the $q$-deformed Dolan-Grady
relations~\cite{Terwilliger2001,Baseilhac2005}:
\begin{equation}
A^3 B - [3]_q A^2 B A + [3]_q A B A^2 - B A^3 = \rho(AB - BA),
\end{equation}
where $\rho = 16$ is the Dolan-Grady coefficient (independent of $q$).
At $q = 1$, $[3]_q = 3$ recovers the classical triple commutator with
binomial coefficients $(1, 3, 3, 1)$. The connection to the Onsager
algebra is: $[2]_1^4 = 2^4 = 16 = \rho$.

%% ══════════════════════════════════════════════════════════════════
\section{$U_q(\mathfrak{sl}_2)$ via \texttt{FreeAlgebra} + \texttt{RingQuot}}

\subsection{Construction}

The quantum group $U_q(\mathfrak{sl}_2)$ is the associative
$\mathbb{k}[q,q^{-1}]$-algebra with generators $E, F, K, K^{-1}$
subject to the Chevalley relations:
\begin{align}
K K^{-1} &= K^{-1} K = 1, \\
K E &= q^2 E K, \quad K F = q^{-2} F K, \\
(q - q^{-1})(EF - FE) &= K - K^{-1}.
\end{align}
The last relation is stated in denominator-free form to avoid
division by $q - q^{-1}$.

In Lean~4:
\begin{enumerate}
\item Define \texttt{Uqsl2Gen} as an inductive type with four constructors.
\item Define \texttt{ChevalleyRel} as an inductive proposition on
\texttt{FreeAlgebra (LaurentPolynomial k) Uqsl2Gen}, encoding all five
relations.
\item Define \texttt{Uqsl2 := RingQuot (ChevalleyRel k)}.
\end{enumerate}
The \texttt{Ring} and \texttt{Algebra (LaurentPolynomial k)} instances
on \texttt{Uqsl2} are \emph{automatic} from \texttt{RingQuot}'s
instance synthesis. This follows the exact pattern of Mathlib's
\texttt{UniversalEnvelopingAlgebra}.

\subsection{Properties}

The quotient map \texttt{uqsl2Mk} sends free algebra elements to
their equivalence classes in $U_q(\mathfrak{sl}_2)$. The generators
$E, F, K, K^{-1}$ are defined as images of the free algebra generators.
$K$ is proved to be a unit (\texttt{IsUnit}) with $K^{-1}$ as its
inverse, using the \texttt{KKinv} and \texttt{KinvK} relations.

All five Chevalley relations are proved as theorems: $KK^{-1} = 1$
and $K^{-1}K = 1$ via \texttt{RingQuot.mkAlgHom\_rel}; $KE = q^2 EK$
and $KF = q^{-2}FK$ via \texttt{AlgHom.commutes}; and the Serre relation
via \texttt{map\_sub} + \texttt{mkAlgHom\_rel} (Aristotle automated prover).

%% ══════════════════════════════════════════════════════════════════
\section{The Chain: Onsager $\to$ Quantum Group $\to$ Gauge Theory}

The formalization of $U_q(\mathfrak{sl}_2)$ is the middle link in a
chain connecting lattice physics to gauge theory:

\begin{center}
\begin{tabular}{lcl}
\textbf{Formalized} & & \textbf{Module} \\
\hline
Onsager algebra $O$ & & OnsagerAlgebra.lean (24 thms) \\
$q$-integers $[n]_q$ & & QNumber.lean (11 thms) \\
$U_q(\mathfrak{sl}_2)$ & & Uqsl2.lean (6 thms + defs) \\
$U_q$ Hopf algebra & & Uqsl2Hopf.lean (21 thms + instances) \\
$D(G)$, $Z(\mathrm{Vec}_G)$ & & 7 Drinfeld/Center modules (96 thms) \\
\hline
\textbf{Planned} & & \\
\hline
$O_q \hookrightarrow U_q(\widehat{\mathfrak{sl}}_2)$ & & coideal embedding \\
$\mathrm{Rep}(U_q(\mathfrak{sl}_2))$ & & MTC at roots of unity \\
\end{tabular}
\end{center}

At $q = 1$, $U_q(\mathfrak{sl}_2)$ reduces to the classical
$U(\mathfrak{sl}_2)$, recovering the Chevalley embedding of the
Onsager algebra. At roots of unity $q = e^{2\pi i/(k+2)}$, the
representation category becomes the $SU(2)_k$ modular tensor
category --- the same structure that governs Chern-Simons gauge theory.

The $q$-Onsager algebra $O_q$ embeds as a coideal subalgebra of the
\emph{affine} quantum group $U_q(\widehat{\mathfrak{sl}}_2)$ (not the
finite type)~\cite{Baseilhac2005,Kolb2014}. This embedding connects
the lattice chirality formalized in companion work to continuous gauge
emergence through the quantum group framework.

%% ══════════════════════════════════════════════════════════════════
\section{Hopf Algebra Structure}

The quantum group $U_q(\mathfrak{sl}_2)$ carries a Hopf algebra
structure: a coproduct $\Delta$, counit $\varepsilon$, and antipode $S$.
We formalize all three and wire them into Mathlib's
\texttt{HopfAlgebra} typeclass, making $U_q(\mathfrak{sl}_2)$ the
\emph{first Hopf algebra in any proof assistant} (beyond the trivial
commutative semiring and group algebra instances).

\subsection{Coproduct and counit}

The coproduct $\Delta : U_q \to U_q \otimes U_q$ is defined on
generators by
\begin{align}
\Delta(E) &= E \otimes K + 1 \otimes E, \\
\Delta(F) &= F \otimes 1 + K^{-1} \otimes F, \\
\Delta(K) &= K \otimes K, \quad \Delta(K^{-1}) = K^{-1} \otimes K^{-1}.
\end{align}
The counit $\varepsilon : U_q \to \mathbb{k}[q,q^{-1}]$ sends
$\varepsilon(E) = \varepsilon(F) = 0$ and $\varepsilon(K) = 1$.

Both are constructed as algebra homomorphisms: we define maps on
generators via \texttt{FreeAlgebra.lift}, prove compatibility with
the Chevalley relations, and descend to the quotient via
\texttt{RingQuot.liftAlgHom}. The \texttt{Bialgebra} instance uses
\texttt{Bialgebra.ofAlgHom}, consuming coassociativity and
counitality proofs.

\subsection{Antipode}

The antipode $S : U_q \to U_q$ is the anti-homomorphism
\begin{equation}
S(E) = -EK^{-1}, \quad S(F) = -KF, \quad S(K) = K^{-1}.
\end{equation}
Since $S$ reverses multiplication ($S(ab) = S(b)S(a)$), we construct
it as an algebra homomorphism into the opposite ring
$U_q^{\mathrm{op}}$ via \texttt{MulOpposite}, then compose with
the canonical linear equivalence. The \texttt{HopfAlgebra} instance
requires the two antipode axioms
$\mu \circ (S \otimes \mathrm{id}) \circ \Delta = \eta \circ \varepsilon$
and
$\mu \circ (\mathrm{id} \otimes S) \circ \Delta = \eta \circ \varepsilon$.

A key structural theorem: the squared antipode is conjugation by $K$,
\begin{equation}
S^2(x) = K x K^{-1},
\end{equation}
which on generators gives $S^2(E) = q^2 E$, $S^2(F) = q^{-2}F$,
$S^2(K) = K$. This inner automorphism (rather than the identity)
reflects the non-cocommutativity of $U_q(\mathfrak{sl}_2)$.

%% ══════════════════════════════════════════════════════════════════
\section{Conclusions}

We have presented the first quantum group in any proof assistant ---
$U_q(\mathfrak{sl}_2)$ as \texttt{RingQuot (ChevalleyRel k)} ---
together with the first non-trivial \texttt{HopfAlgebra} instance,
all with zero axioms. The construction uses Mathlib's
\texttt{FreeAlgebra} $\to$ \texttt{RingQuot} $\to$
\texttt{liftAlgHom} pattern for the algebra, and the
\texttt{Coalgebra} $\to$ \texttt{Bialgebra} $\to$
\texttt{HopfAlgebra} hierarchy for the coalgebra structure.

\begin{acknowledgments}
Lean~4 proofs verified by \texttt{lake build} (v4.28.0, Mathlib \texttt{8f9d9cff}).
Automated proofs by the Aristotle theorem prover (Harmonic).
\end{acknowledgments}

\begin{thebibliography}{99}
\bibitem{Drinfeld1986}
V.~G.~Drinfeld, Proc.\ ICM Berkeley, 798 (1986).

\bibitem{Jimbo1986}
M.~Jimbo, Lett.\ Math.\ Phys.\ \textbf{11}, 247 (1986).

\bibitem{Terwilliger2001}
P.~Terwilliger, arXiv:math/0307016.

\bibitem{Baseilhac2005}
P.~Baseilhac and K.~Belliard,
Nucl.\ Phys.\ B \textbf{873}, 550 (2013) [arXiv:0906.1215].

\bibitem{Kolb2014}
S.~Kolb, Adv.\ Math.\ \textbf{267}, 395 (2014) [arXiv:1206.2117].

\bibitem{Kassel1995}
C.~Kassel, \emph{Quantum Groups} (Springer, 1995).

\bibitem{Paper9}
J.~Roehm, companion paper: SM anomaly in $\mathbb{Z}_{16}$ + Drinfeld center.

\bibitem{Paper10}
J.~Roehm, companion paper: From modular forms to generation counting.
\end{thebibliography}

\end{document}
